% problem-000151 1 晶体结构与晶胞
\section*{1. 晶体结构与晶胞}
\textit{下列为分问。}

\subsection*{1-1}
以阳极泥中的 $\mathrm { A g _ { 2 } S e }$ 和 $\mathrm { C u _ { 2 } S e }$ 为例，写出反应1中的化学反应⽅程式。

\subsection*{1-2}
写出反应2中的化学反应⽅程式。

\subsection*{1-3}
双β衰变是原⼦核的⼀种稀有衰变⽅式，可近似看作连续的两次β衰变。第⼀个成功观测到双β衰变现象的原⼦核是 $^ { 8 2 } \mathrm { S e }$ ，写出82Se发⽣双β衰变的反应⽅程式。

\subsection*{1-4}
$\mathrm { G a C l _ { 3 } }$ 和 $\mathrm { S e C l _ { 4 } }$ 可在室温下发⽣反应，⽣成离⼦型化合物R。然⽽，将等摩尔数的 $S e O _ { 2 } ,$ $\mathrm { S e C l _ { 4 } }$ 和 $\mathrm { G a C l _ { 3 } }$ 于 $5 0 ~ ^ { \circ } \mathrm { C }$ 反应，得到等摩尔数的⽆⾊晶体P和液态物质Q。单晶X-射线衍射测试表明，⽆⾊晶体P中，Ga中⼼为四⾯体配位⼏何构型，Se的配位⼏何为三⻆锥形。Q的分⼦⼏何构型为三⻆锥形。在 $8 0 0 ~ ^ { \circ } \mathrm { C }$ 惰性⽓氛中，P发⽣热分解，剩余残渣的质量百分含量约为 $1 1 \text{‰}$

\subsection*{1-4}
-1 写出能表明R结构特征的化学式，分别写出P和Q的化学式。

\subsection*{1-4}
-2 根据Se、Ga氯化物和氧化物的相关性质，推导出P的热分解产物，并写出化学反应⽅程式。
